% Môn: Toán
% Lớp: 10
% Chương: Chương I. Mệnh đề và tập hợp
% Bài:  Mệnh đề · Mệnh đề kéo theo
% Số câu: 0
% App đọc file này trực tiếp — sửa rồi Commit trên GitHub
% Ghi nguồn từng câu: \nguon{SGK} hoặc \nguon{2 SBT VL 12 KNTT} trong


%%%=============EX_1=============%%%

% Mức: TH
% ID: T10C1B1-139-TN
% ID: T10C1B1-139-TN
% Mức: TH
% ID: T10C1B1-139-TN

% Mức: TH
% ID: T10C1B1-140-TN
% ID: T10C1B1-140-TN
% Mức: TH
% ID: T10C1B1-140-TN

% Mức: TH
% ID: T10C1B1-141-TN
% ID: T10C1B1-141-TN
% Mức: TH
% ID: T10C1B1-141-TN

% Mức: TH
% ID: T10C1B1-142-TN
% ID: T10C1B1-142-TN
% Mức: TH
% ID: T10C1B1-142-TN

% Mức: TH
% ID: T10C1B1-143-TN
% ID: T10C1B1-143-TN
% Mức: TH
% ID: T10C1B1-143-TN

% Mức: TH
% ID: T10C1B1-144-TN
% ID: T10C1B1-144-TN
% Mức: TH
% ID: T10C1B1-144-TN

% Mức: TH
% ID: T10C1B1-145-TN
% ID: T10C1B1-145-TN
% Mức: TH
% ID: T10C1B1-145-TN

% Mức: NB
% ID: T10C1B1-146-TN
% ID: T10C1B1-146-TN
% Mức: NB
% ID: T10C1B1-146-TN

%%%=============EX_15=============%%%

% Mức: NB
% ID: T10C1B1-147-TN
% ID: T10C1B1-147-TN
% Mức: NB
% ID: T10C1B1-147-TN

%%%=============EX_16=============%%%

% Mức: NB
% ID: T10C1B1-148-TN
% ID: T10C1B1-148-TN
% Mức: NB
% ID: T10C1B1-148-TN

%%%=============EX_17=============%%%

% Mức: NB
% ID: T10C1B1-149-TN
% ID: T10C1B1-149-TN
% Mức: NB
% ID: T10C1B1-149-TN

%%%=============EX_18=============%%%

% Mức: NB
% ID: T10C1B1-150-TN
% ID: T10C1B1-150-TN
% Mức: NB
% ID: T10C1B1-150-TN
% Mức: NB

%%%=============EX_19=============%%%
\dangbt{Mệnh đề kéo theo}
\begin{ex}
	% ID: T10C1B1-150-TN
	% Mức: NB
	Trong các mệnh đề kéo theo sau, mệnh đề nào là mệnh đề đúng\
	\choice
		{ Nếu ${37}$ là số nguyên thì ${11}$ là hợp số  }
		{ \True Nếu ${38}$ là số chẵn thì ${60}$ không phải là số nguyên tố }
		{ Nếu ${45}$ là số lẻ thì ${46}$ là số lẻ }
		{ Nếu ${13}$ là số nguyên tố thì ${6}$ là số lẻ }
	\loigiai{
		Nếu ${38}$ là số chẵn thì ${60}$ không phải là số nguyên tố là khẳng định đúng
		}
\end{ex}

\dangbt{Mệnh đề kéo theo}
\begin{ex}
	Cho định lí: "Nếu hai tam giác bằng nhau thì diện tích của chúng bằng nhau". Mệnh đề nào sau đây là phát biểu đúng về định lí này?
	\choice
	{Hai tam giác bằng nhau là điều kiện cần để diện tích của chúng bằng nhau.}
	{Hai tam giác có diện tích bằng nhau là điều kiện đủ để chúng bằng nhau.}
	{\True Hai tam giác bằng nhau là điều kiện đủ để diện tích của chúng bằng nhau.}
	{Hai tam giác bằng nhau là điều kiện cần và đủ để diện tích của chúng bằng nhau.}
	\loigiai{
		Với định lí dạng $P \Rightarrow Q$, ta có cách phát biểu: "$P$ là điều kiện đủ để có $Q$" hoặc "$Q$ là điều kiện cần để có $P$". Ở đây $P$ là "Hai tam giác bằng nhau" và $Q$ là "Diện tích của chúng bằng nhau". Do đó, "Hai tam giác bằng nhau là điều kiện đủ để diện tích của chúng bằng nhau" là phát biểu đúng.
	}
	\nguon{SGK KNTT Toán Lớp 10 Bài 1 trang 8}
\end{ex}

\dangbt{Mệnh đề kéo theo}
\begin{ex}
	Cho mệnh đề $P$: "Nếu $a+b < 2$ thì một trong hai số $a$ và $b$ nhỏ hơn $1$". Mệnh đề nào sau đây là một cách phát biểu khác của mệnh đề đã cho?
	\choice
	{\True Điều kiện đủ để một trong hai số $a$ và $b$ nhỏ hơn $1$ là $a+b < 2$.}
	{Điều kiện cần để một trong hai số $a$ và $b$ nhỏ hơn $1$ là $a+b < 2$.}
	{Điều kiện đủ để $a+b < 2$ là một trong hai số $a$ và $b$ nhỏ hơn $1$.}
	{Điều kiện cần và đủ để một trong hai số $a$ và $b$ nhỏ hơn $1$ là $a+b < 2$.}
	\loigiai{
		Mệnh đề kéo theo $P \Rightarrow Q$ có dạng: "Nếu $P$ thì $Q$". Ta phát biểu lại là "$P$ là điều kiện đủ để có $Q$" hay "Điều kiện đủ để có $Q$ là $P$". Với $P$ là "$a+b < 2$" và $Q$ là "một trong hai số $a$ và $b$ nhỏ hơn $1$". Vậy cách phát biểu khác là: "Điều kiện đủ để một trong hai số $a$ và $b$ nhỏ hơn $1$ là $a+b < 2$".
	}
	\nguon{SBT KNTT Toán Lớp 10 Bài 1 Câu 1.18}
\end{ex}

\dangbt{Mệnh đề kéo theo}
\begin{ex}
	Trong các mệnh đề kéo theo sau, mệnh đề nào là mệnh đề đúng?
	\choice
	{Nếu "33 là hợp số" thì "15 chia hết cho 25".}
	{Nếu "7 là số nguyên tố" thì "8 là bội số của 3".}
	{\True Nếu "20 là hợp số" thì "24 chia hết cho 6".}
	{Nếu "3 + 9 = 12" thì "4 > 7".}
	\loigiai{
		Mệnh đề kéo theo $P \Rightarrow Q$ chỉ sai khi $P$ đúng và $Q$ sai.
		- Xét phương án C: Mệnh đề giả thiết "20 là hợp số" là đúng; mệnh đề kết luận "24 chia hết cho 6" là đúng. Vì cả $P$ và $Q$ đều đúng nên mệnh đề kéo theo $P \Rightarrow Q$ là đúng.
		- Ở các phương án khác, giả thiết đúng nhưng kết luận sai, dẫn tới mệnh đề kéo theo là sai.
	}
	\nguon{Toán 10 Học kì 1 GV.pdf Bài 1 Câu 19}
\end{ex}

\dangbt{Mệnh đề kéo theo}
\begin{ex}
	Định lí có dạng $A \Rightarrow B$ được phát biểu theo ngôn ngữ điều kiện cần và điều kiện đủ như thế nào là đúng?
	\choice
	{$A$ là điều kiện cần để có $B$.}
	{\True $A$ là điều kiện đủ để có $B$.}
	{$B$ là điều kiện đủ để có $A$.}
	{$A$ là điều kiện cần và đủ để có $B$.}
	\loigiai{
		Với định lí có dạng $A \Rightarrow B$, ta phát biểu:
		- $A$ là giả thiết, $B$ là kết luận.
		- $A$ là điều kiện đủ để có $B$.
		- $B$ là điều kiện cần để có $A$.
		Do đó phương án B là khẳng định đúng.
	}
	\nguon{Toán 10 Học kì 1 GV.pdf Bài 1 Câu 9}
\end{ex}

\dangbt{Mệnh đề kéo theo}
\begin{ex}
	Cho định lí: "Nếu một số tự nhiên có tận cùng là $0$ thì số đó chia hết cho $10$". Phát biểu nào sau đây sử dụng khái niệm "điều kiện cần" để mô tả định lí trên?
	\choice
	{Một số tự nhiên có tận cùng là $0$ là điều kiện cần để số đó chia hết cho $10$.}
	{\True Một số tự nhiên chia hết cho $10$ là điều kiện cần để số đó có tận cùng là $0$.}
	{Một số tự nhiên chia hết cho $10$ là điều kiện đủ để số đó có tận cùng là $0$.}
	{Một số tự nhiên có tận cùng là $0$ là điều kiện cần và đủ để số đó chia hết cho $10$.}
	\loigiai{
		Với định lí có dạng $P \Rightarrow Q$: "Nếu một số tự nhiên có tận cùng là $0$ ($P$) thì số đó chia hết cho $10$ ($Q$)", ta phát biểu "$Q$ là điều kiện cần để có $P$". Vậy: "Một số tự nhiên chia hết cho $10$ là điều kiện cần để số đó có tận cùng là $0$".
	}
	\nguon{SBT KNTT Toán Lớp 10 Bài 1 Câu 1.33a}
\end{ex}

\dangbt{Mệnh đề kéo theo}
\begin{ex}
	Cho hai mệnh đề $P$: "Tứ giác $ABCD$ là hình bình hành" và $Q$: "Tứ giác $ABCD$ có hai cạnh đối song song và bằng nhau". Phát biểu nào sau đây đúng về mệnh đề kéo theo $P \Rightarrow Q$?
	\choice
	{Nếu tứ giác $ABCD$ có hai cạnh đối song song và bằng nhau thì tứ giác $ABCD$ là hình bình hành.}
	{Tứ giác $ABCD$ là hình bình hành kéo theo tứ giác $ABCD$ có hai đường chéo cắt nhau tại trung điểm mỗi đường.}
	{\True Nếu tứ giác $ABCD$ là hình bình hành thì tứ giác $ABCD$ có hai cạnh đối song song và bằng nhau.}
	{Tứ giác $ABCD$ có hai cạnh đối song song và bằng nhau là điều kiện đủ để tứ giác $ABCD$ là hình bình hành.}
	\loigiai{
		Mệnh đề kéo theo $P \Rightarrow Q$ phát biểu dưới dạng "Nếu $P$ thì $Q$".
		Với $P$: "Tứ giác $ABCD$ là hình bình hành" và $Q$: "Tứ giác $ABCD$ có hai cạnh đối song song và bằng nhau", mệnh đề kéo theo $P \Rightarrow Q$ được phát biểu là: "Nếu tứ giác $ABCD$ là hình bình hành thì tứ giác $ABCD$ có hai cạnh đối song song và bằng nhau".
	}
	\nguon{SGK KNTT Toán Lớp 10 Bài 1 trang 8}
\end{ex}

\dangbt{Mệnh đề kéo theo}
\begin{ex}
	Cho mệnh đề $A$: "Với mọi số thực $x$, nếu $x < -2$ thì $x^2 > 4$". Khẳng định nào sau đây đúng?
	\choice
	{\True Mệnh đề $A$ đúng vì với mọi số thực $x < -2$ thì bình phương của nó luôn lớn hơn $4$.}
	{Mệnh đề $A$ sai vì với $x = -3$ ta có $x^2 = 9 > 4$ là sai.}
	{Mệnh đề $A$ sai vì với $x = -1 < -2$ ta có $x^2 = 1 > 4$ là sai.}
	{Mệnh đề $A$ đúng vì với $x = -2$ ta có $x^2 = 4 > 4$ là đúng.}
	\loigiai{
		Mệnh đề có dạng: "Với mọi số thực $x$, nếu $x < -2$ thì $x^2 > 4$".
		Khi $x < -2 \Rightarrow -x > 2 \Rightarrow (-x)^2 > 2^2 \Rightarrow x^2 > 4$. Khẳng định này hoàn toàn đúng với mọi số thực $x$. Do đó mệnh đề $A$ là mệnh đề đúng.
	}
	\nguon{SBT KNTT Toán Lớp 10 Bài 1 Câu 1.31A}
\end{ex}

\dangbt{Mệnh đề kéo theo}
\begin{ex}
	Cho mệnh đề kéo theo: "Nếu $x^2 < 4$ thì $x < -2$". Xét tính đúng sai của mệnh đề này và tìm một phản ví dụ (nếu có).
	\choice
	{Mệnh đề đúng.}
	{Mệnh đề sai, phản ví dụ là $x = -3$.}
	{\True Mệnh đề sai, phản ví dụ là $x = 0$.}
	{Mệnh đề sai, phản ví dụ là $x = 3$.}
	\loigiai{
		Mệnh đề đã cho có dạng $P \Rightarrow Q$. Xét $x = 0$, ta có $0^2 = 0 < 4$ (giả thiết đúng) nhưng $0 < -2$ là sai (kết luận sai). Do đó mệnh đề $P \Rightarrow Q$ là mệnh đề sai, và $x = 0$ là một phản ví dụ.
	}
	\nguon{SBT KNTT Toán Lớp 10 Bài 1 Câu 1.31B}
\end{ex}

\dangbt{Mệnh đề kéo theo}
\begin{ex}
	Cho hai mệnh đề $P$: "Hai tam giác bằng nhau" và $Q$: "Hai tam giác có diện tích bằng nhau". Xét tính đúng sai của các khẳng định sau:
	\choiceTF
	{\True Mệnh đề $P \Rightarrow Q$ là: "Nếu hai tam giác bằng nhau thì chúng có diện tích bằng nhau".}
	{\True Mệnh đề $Q \Rightarrow P$ là: "Nếu hai tam giác có diện tích bằng nhau thì hai tam giác đó bằng nhau".}
	{Mệnh đề đảo $Q \Rightarrow P$ là một mệnh đề đúng.}
	{\True Mệnh đề kéo theo $P \Rightarrow Q$ là một mệnh đề đúng.}
	\loigiai{
		- Ý thứ nhất Đúng: Đây là phát biểu chuẩn xác của mệnh đề kéo theo $P \Rightarrow Q$.
		- Ý thứ hai Đúng: Đây là phát biểu chuẩn xác của mệnh đề đảo $Q \Rightarrow P$.
		- Ý thứ ba Sai: Hai tam giác có diện tích bằng nhau chưa chắc đã bằng nhau (ví dụ một tam giác có cạnh đáy bằng $4$, chiều cao bằng $3$ và một tam giác có cạnh đáy bằng $6$, chiều cao bằng $2$ đều có diện tích bằng $6$ nhưng chúng không bằng nhau).
		- Ý thứ tư Đúng: Hai tam giác bằng nhau thì chắc chắn diện tích của chúng phải bằng nhau.
	}
	\nguon{Toán 10 Học kì 1 GV.pdf Bài 1 Câu 4}
\end{ex}

\dangbt{Mệnh đề kéo theo}
\begin{ex}
	Cho hai mệnh đề $A$: "$a + b + c = 0$" và $B$: "Phương trình bậc hai $ax^2 + bx + c = 0$ có một nghiệm bằng $1$". Xét tính đúng sai của các khẳng định sau:
	\choiceTF
	{\True Mệnh đề kéo theo $A \Rightarrow B$ là một mệnh đề đúng.}
	{\True Mệnh đề đảo $B \Rightarrow A$ là một mệnh đề đúng.}
	{Biểu thức $a + b + c \neq 0$ là điều kiện đủ để phương trình $ax^2 + bx + c = 0$ có một nghiệm bằng $1$.}
	{\True Ta có thể phát biểu: "Điều kiện cần và đủ để phương trình $ax^2 + bx + c = 0$ có một nghiệm bằng $1$ là $a + b + c = 0$".}
	\loigiai{
		- Ý thứ nhất Đúng: Nếu $a + b + c = 0$, ta thế $x = 1$ vào phương trình $ax^2 + bx + c = 0$ ta được $a \cdot 1^2 + b \cdot 1 + c = a + b + c = 0$ (đúng). Nên phương trình luôn có nghiệm $x = 1$.
		- Ý thứ hai Đúng: Nếu phương trình $ax^2 + bx + c = 0$ có một nghiệm bằng $1$ thì khi thế $x = 1$ vào phương trình ta phải được đẳng thức đúng, tức là $a \cdot 1^2 + b \cdot 1 + c = 0 \Leftrightarrow a + b + c = 0$.
		- Ý thứ ba Sai: Vì biểu thức $a + b + c = 0$ mới là điều kiện đủ, còn nếu $a + b + c \neq 0$ thì phương trình không thể nhận $1$ làm nghiệm.
		- Ý thứ tư Đúng: Do hai mệnh đề kéo theo thuận nghịch đều đúng nên ta có điều kiện cần và đủ.
	}
	\nguon{Toán 10 Học kì 1 GV.pdf Bài 1 Bài 8}
\end{ex}

\dangbt{Mệnh đề kéo theo}
\begin{ex}
	Xét tính đúng sai của mệnh đề kéo theo $P \Rightarrow Q$: "Vì $120$ chia hết cho $6$ nên $120$ chia hết cho $9$". Ghi "Đúng" nếu mệnh đề đúng, ghi "Sai" nếu mệnh đề sai.
	\shortans{Sai}
	\loigiai{
		Mệnh đề kéo theo $P \Rightarrow Q$ chỉ sai khi $P$ đúng và $Q$ sai.
		Ở đây, giả thiết $P$: "120 chia hết cho 6" là đúng (vì $120 = 6 \cdot 20$); nhưng kết luận $Q$: "120 chia hết cho 9" là sai (vì tổng các chữ số là $1+2+0=3$ không chia hết cho $9$). Vì $P$ đúng và $Q$ sai nên mệnh đề kéo theo đã cho là một mệnh đề sai.
	}
	\nguon{Toán 10 Học kì 1 GV.pdf Bài 1 Bài 3a}
\end{ex}

\dangbt{Mệnh đề kéo theo}
\begin{ex}
	Cho mệnh đề kéo theo $A$: "Nếu $120$ chia hết cho $6$ thì $120$ chia hết cho $9$". Xét tính đúng sai của mệnh đề đảo của mệnh đề $A$. Ghi "Đúng" nếu đúng, ghi "Sai" nếu sai.
	\shortans{Đúng}
	\loigiai{
		Mệnh đề đảo của mệnh đề $A$ là: "Nếu $120$ chia hết cho $9$ thì $120$ chia hết cho $6$".
		Xét tính đúng sai: Giả thiết "120 chia hết cho 9" là một khẳng định sai. Trong logic toán học, một mệnh đề kéo theo có giả thiết sai thì luôn luôn đúng (bất kể kết luận đúng hay sai). Do đó, mệnh đề đảo này là một mệnh đề đúng.
	}
	\nguon{Toán 10 Học kì 1 GV.pdf Bài 1 Bài 3b}
\end{ex}

\dangbt{Mệnh đề kéo theo}
\begin{ex}
	Cho định lí: "Nếu tứ giác $ABCD$ là hình bình hành thì tứ giác $ABCD$ có hai đường chéo cắt nhau tại trung điểm mỗi đường". Xét tính đúng sai của mệnh đề đảo của định lí trên. Ghi "Đúng" nếu đúng, ghi "Sai" nếu sai.
	\shortans{Đúng}
	\loigiai{
		- Mệnh đề đảo của định lí là: "Nếu tứ giác $ABCD$ có hai đường chéo cắt nhau tại trung điểm mỗi đường thì tứ giác $ABCD$ là hình bình hành".
		- Theo hình học phẳng, đây chính là một dấu hiệu nhận biết hình bình hành đã được chứng minh là đúng. Do đó mệnh đề đảo này là một mệnh đề đúng.
	}
	\nguon{Toán 10 Học kì 1 GV.pdf Bài 1 Bài 4b}
\end{ex}

\dangbt{Mệnh đề kéo theo}
\begin{ex}
	Cho hai mệnh đề $P$: "$a$ và $b$ cùng chia hết cho $c$" và $Q$: "$a + b$ chia hết cho $c$" (với $a, b, c$ là các số nguyên, $c \ne 0$).
	a) Hãy phát biểu định lí $P \Rightarrow Q$. Nêu giả thiết và kết luận của định lí.
	b) Phát biểu định lí này dưới dạng "điều kiện cần" và "điều kiện đủ".
	\loigiai{
		a) Phát biểu định lí $P \Rightarrow Q$: "Nếu $a$ và $b$ cùng chia hết cho $c$ thì $a + b$ chia hết cho $c$".
		- Giả thiết của định lí là: "$a$ và $b$ cùng chia hết cho $c$".
		- Kết luận của định lí là: "$a + b$ chia hết cho $c$".
		b) Phát biểu dưới dạng điều kiện cần và điều kiện đủ:
		- Sử dụng "điều kiện đủ": "$a$ và $b$ cùng chia hết cho $c$ là điều kiện đủ để $a + b$ chia hết cho $c$".
		- Sử dụng "điều kiện cần": "$a + b$ chia hết cho $c$ là điều kiện cần để $a$ và $b$ cùng chia hết cho $c$".
	}
	\nguon{Toán 10 Học kì 1 GV.pdf Bài 1 Ví dụ 4a}
\end{ex}

\dangbt{Mệnh đề kéo theo}
\begin{ex}
	Hãy phát biểu mệnh đề đảo của mệnh đề sau và xác định tính đúng sai của mệnh đề đảo đó: "Nếu tam giác $ABC$ là tam giác đều thì tam giác $ABC$ là tam giác cân".
	\loigiai{
		- Mệnh đề đảo của mệnh đề đã cho là: "Nếu tam giác $ABC$ là tam giác cân thì tam giác $ABC$ là tam giác đều".
		- Xét tính đúng sai: Mệnh đề đảo này là một mệnh đề sai. Vì một tam giác cân chưa chắc đã là tam giác đều (chẳng hạn tam giác cân có một góc ở đỉnh bằng $50^\circ$ thì không thể là tam giác đều).
	}
	\nguon{Toán 10 Học kì 1 GV.pdf Bài 1 Bài 4}
\end{ex}

\dangbt{Mệnh đề kéo theo}
\begin{ex}
	Phát biểu mệnh đề đảo của mỗi mệnh đề sau và xác định tính đúng sai của chúng:
	a) $P$: "Nếu số tự nhiên $n$ có chữ số tận cùng là $5$ thì $n$ chia hết cho $5$";
	b) $Q$: "Nếu tứ giác $ABCD$ là hình chữ nhật thì tứ giác $ABCD$ có hai đường chéo bằng nhau".
	\loigiai{
		a) Mệnh đề đảo của $P$: "Nếu số tự nhiên $n$ chia hết cho $5$ thì $n$ có chữ số tận cùng là $5$".
		- Mệnh đề đảo này là một mệnh đề sai. Ví dụ phản chứng: Số tự nhiên $n = 10$ chia hết cho $5$ nhưng có chữ số tận cùng bằng $0$ chứ không phải bằng $5$.
		b) Mệnh đề đảo của $Q$: "Nếu tứ giác $ABCD$ có hai đường chéo bằng nhau thì tứ giác $ABCD$ là hình chữ nhật".
		- Mệnh đề đảo này là một mệnh đề sai. Ví dụ phản chứng: Tứ giác $ABCD$ là hình thang cân thì nó có hai đường chéo bằng nhau nhưng không phải là hình chữ nhật.
	}
	\nguon{SGK KNTT Toán Lớp 10 Bài 1 trang 11 Bài 1.4}
\end{ex}

\dangbt{Mệnh đề kéo theo}
\begin{ex}
	Với hai số thực $a$ và $b$, xét các mệnh đề $P$: "$a^2 < b^2$" và $Q$: "$0 < a < b$".
	a) Hãy phát biểu mệnh đề kéo theo $P \Rightarrow Q$ và phát biểu mệnh đề đảo của nó.
	b) Xác định tính đúng sai của mỗi mệnh đề đó.
	\loigiai{
		a) Phát biểu các mệnh đề:
		- Mệnh đề kéo theo $P \Rightarrow Q$: "Nếu $a^2 < b^2$ thì $0 < a < b$".
		- Mệnh đề đảo $Q \Rightarrow P$: "Nếu $0 < a < b$ thì $a^2 < b^2$".
		b) Xét tính đúng sai:
		- Mệnh đề kéo theo $P \Rightarrow Q$ là một mệnh đề sai. Ví dụ phản chứng: Chọn $a = -1$ và $b = 2$, ta có $a^2 = (-1)^2 = 1 < b^2 = 2^2 = 4$ (giả thiết đúng) nhưng không có $0 < -1 < 2$ vì $-1 < 0$ (kết luận sai).
		- Mệnh đề đảo Q \Rightarrow P$ là một mệnh đề đúng vì từ $0 < a < b$, nhân cả hai vế của bất đẳng thức $a < b$ lần lượt với các số thực dương $a$ và $b$, ta được $a^2 < ab$ và $ab < b^2$, từ đó suy ra $a^2 < b^2.
	}
	\nguon{SGK KNTT Toán Lớp 10 Bài 1 trang 11 Bài 1.5}
\end{ex}